% Quelle: 3. Potenzen in Zm.pdf
% Art: Vorlesung
% Themen: schnelles Potenzieren mod m, Binärdarstellung des Exponenten, Einheitengruppe Zm*, Kürzungsregel, Satz von Fermat, Ordnung eines Elements, zyklische Gruppe, primitives Element, Satz von Euler, Eulersche phi-Funktion, RSA-Verschlüsselung, öffentlicher und privater Schlüssel

\section{Vorlesung 3: Potenzen in $\Zm{m}$}
\subsection{3.1 Schnelle Berechnung von $a^k \pmod{m}$}

Man schreibt im Restklassenring $\Zm{m}$, wie im Körper $\R$
\[
  a^k := a \cdot a \cdot \ldots \cdot a \qquad (a \neq 0)
\]
Wie immer, wenn es um Multiplikation geht, schließen wir die 0 aus.

\begin{algorithmus}[Schnelles Potenzieren]
Nutzen die Binärdarstellung von
\[
  k = k_0 \cdot 2^0 + k_1 \cdot 2^1 + k_2 \cdot 2^2 + \ldots + k_r \cdot 2^r . \qquad k_i \in \{0, 1\}
\]
\end{algorithmus}

\begin{beispiel}
Gesucht $12^{100} \pmod{34}$
\begin{align*}
  & 100 = 2^6 + 2^5 + 2^2 \\
  \Rightarrow\quad & 12^{100} = 12^{64} \cdot 12^{32} \cdot 12^{4} \\
  & 12 \\
  & 12^2 \equiv 136 + 8 \pmod{34} \\
  & 12^4 \equiv 8^2 \equiv -4 \pmod{34} \\
  & 12^8 \equiv (-4)^2 = 16 \pmod{34} \\
  & 12^{16} \equiv 16^2 = 272 - 16 \equiv -16 \pmod{34} \\
  & 12^{32} \equiv (-16)^2 \equiv -16 \pmod{34} \\
  & 12^{64} \equiv (-16)^2 \equiv -16 \pmod{34}
\end{align*}
Also
\[
  12^{100} \equiv (-16) \cdot (-16) \cdot (-4) \equiv (-16) \cdot (-4) \equiv 64 \equiv -4 \equiv 30 \pmod{34}
\]
\end{beispiel}

\subsection{3.2 Die Einheitengruppe $\Zm{m}^*$ von $\Zm{m}$}

\begin{bemerkung}
Die Einheiten sind genau die Elemente $k \in \Zm{m}$ mit $\ggT(k, m) = 1$.\\
$(\Zm{m}^*, \cdot\,)$ ist eine abelsche Gruppe,
\[
  \operatorname{ord}(G) = |G|, \qquad\qquad \operatorname{ord}(\Zm{m}^*) = \varphi(m).
\]
\end{bemerkung}

\begin{bemerkung}[Kürzungsregel]
Für Einheiten $a \in \Zm{m}^*$ gilt
\[
  a \cdot x = a \cdot y \quad \Rightarrow \quad x = y \qquad\qquad \text{Kürzungsregel}
\]
\end{bemerkung}

Sei $p$ eine Primzahl. Dann ist
\begin{align*}
  \Rightarrow\quad & \Zm{p}^* = \{1, 2, 3, \ldots, p - 1\} = \{a1, a2, a3, \ldots, a(p - 1)\} \\
  \intertext{Dann stimmen auch die Produkte überein}
  \Rightarrow\quad & 1 \cdot 2 \cdot 3 \cdot \ldots \cdot (p - 1) = a^{p-1} \cdot 1 \cdot 2 \cdot 3 \cdot \ldots \cdot (p - 1) \\
  \Rightarrow\quad & 1 = a^{p-1}
\end{align*}

\begin{beispiel}
$p = 7$, $a = 3$ $\Rightarrow$ $\{a1, a2, a3, \ldots, a6\} = \{3, 6, 2, 5, 1, 4\}$
\end{beispiel}

\begin{satz}[Fermat]
Sei $p$ eine Primzahl und $a \in \Zm{p}^*$, dann gilt \hfill $a^{\operatorname{ord} G} = 1$.
\begin{align*}
  & a^{p-1} \equiv 1 \pmod{p} \\
  \text{bzw.}\quad & a^{p-1} = 1 \quad (\text{in } \Zm{p}). \qquad \square
\end{align*}
\end{satz}

Im Beweis haben wir genutzt, dass $\Zm{p}^*$ abelsch ist.

\begin{definition}[Bezeichnung]
Die kleinste Zahl $k > 0$ mit $a^k = 1$ heißt \textbf{Ordnung von $a$} in $G$. Die Menge
\[
  \langle a \rangle := \{\, a, a^2, a^3, \ldots, a^k = 1 \,\} \subseteq G
\]
ist eine zyklische Gruppe der Ordnung $k$.

Ein Element $g \in G$ mit $\operatorname{ord}(g) = |G|$ heißt erzeugendes Element (oder \textbf{primitiv}) von $G$. Für solches $g$ gilt
\[
  G = \{\, g, g^2, \ldots, g^k = 1 \} = \langle g \rangle
\]
Besitzt $G$ ein primitives Element, so heißt $G$ zyklisch. Zyklische Gruppen sind abelsch.
\end{definition}

\begin{beispiel}
$\Zm{7}^* = \{1, 2, \ldots, 6\}$
\begin{align*}
  & 2, 2^2, 2^3, \ldots, 2^6 = 2, 4, 1, 2, 4, 1 && \Rightarrow\ \operatorname{ord}(2) = 3 \\
  & 3, 3^2, 3^3, \ldots, 3^6 = 3, 2, -1, -3, -2, 1 && \operatorname{ord}(3) = 6 \\
  \Rightarrow\quad & 3 \text{ ist primitives Element von } \Zm{7}^*.
\end{align*}
\end{beispiel}

\begin{bemerkung}
Fermat war begeistert von seinem Satz. Er zeige, dass die Griechen noch nicht alles wußten.

Euler hat den Satz von Fermat etwa 100 Jahre später auf $\Zm{m}^*$ übertragen.\\
$\operatorname{ord} \Zm{m}^* = \boldsymbol{\varphi}(m)$.
\end{bemerkung}

\begin{satz}[Euler]
Formulierung in $\Zm{m}$
\[
  a^{\varphi(m)} = 1 \qquad\qquad \text{für } a \in \Zm{m}^*
\]
Formulierung in $\Z$
\[
  a^{\varphi(m)} \equiv 1 \pmod{m} \qquad \text{für } \ggT(a, m) = 1
\]
\end{satz}

\begin{beispiel}
Fortsetzung: $\varphi(34) = \varphi(2 \cdot 17) = \varphi(2) \cdot \varphi(17) = 1 \cdot 16$
\[
  \Rightarrow \quad 12^{16} = -16 \neq 1 \pmod{34}
\]
Kein Wunder: $\ggT(12, 34) = 2$ teilt jede Potenz von 12 mod 34.
\end{beispiel}

\begin{bemerkung}
Falls $\ggT(a, m) = 1$, geht das Potenzieren schneller:
\begin{align*}
  & k = q \cdot \varphi(m) + r \equiv r && (\bmod\ \varphi(m)) \\
  \Rightarrow\quad & a^k = (a^{\varphi(m)})^q \cdot a^r \equiv 1^q \cdot a^r = a^r && (\bmod\ m) \qquad \square
\end{align*}
\end{bemerkung}

\begin{beispiel}
Gesucht: $11^{16} \pmod{34}$
\begin{align*}
  & 11 \\
  & 11^2 = 121 = 3 \cdot 34 + 19 \\
  & 11^4 \equiv 19^2 = 361 = 10 \cdot 34 + 21 \pmod{34} \\
  & 11^8 \equiv 21^2 = 441 = 13 \cdot 34 - 1 \pmod{34} \\
  & 11^{16} \equiv (-1)^2 = 1 \pmod{34}
\end{align*}
\end{beispiel}

\subsection{3.3 Die RSA-Verschlüsselung}

Ändern nun die Blickrichtung und fragen, für welche Exponenten $e\,d$ ist
\[
  a^{e\,d} \equiv a \qquad (\bmod\ m)?
\]
Multiplizieren die Euler-Formel mit $a$ und finden
\[
  a^{1 + \varphi(m)} \equiv a \pmod{m}
\]
also \quad $e\,d = 1$ oder $1 + \boldsymbol{\varphi}(m)$ \quad allgemein: \quad $e\,d = 1 \pmod{\varphi(m)}$

\begin{definition}[RSA: Ver- und Entschlüsselung]
So kann man das Klartextwort $a \in \Z$ verschlüsseln: \quad $a \to c = a^e$\\
$c \in \Z$ ist das Codewort

entschlüsseln: \quad $a \to c = a^e \to a = a^{ed} \pmod{m}$

Man nennt $e$ den Verschlüsselungsexponenten (encoding exponent) und\\
$d$ den Entschlüsselungsexponenten (decoding exponent).\\
Die einzige Bedingung an $e$ ist \ $e \in \Zm{\varphi(m)}^*$, also $\ggT(e, \varphi(m)) = 1$.\\
$d$ ist die Inverse von $e$ in $\Zm{\varphi(m)}^*$.
\end{definition}

\begin{beispiel}
Sei $m = 101$, $e = 13 \in \Zm{\varphi(m)}^*$, \quad $\varphi(101) = 100$.\\
$13 \cdot 7 \cdot 11 = 1001 \equiv 1 \pmod{100}$, also\\
$d = 77$. Verschlüsselt werden soll die Zahl $a = 8 \in \Zm{m}^*$
\[
  c \equiv a^e = 8^{13} \equiv 18 \qquad (\bmod\ 101) \qquad (\text{Codewort})
\]
\begin{align*}
  & 13 = 1 + 2^2 + 2^3, \\
  & 8^{13} = 8^8 \cdot 8^4 \cdot 8 \\
  & 8^2 = 64 \\
  & 8^4 = 4096 \equiv 56 \pmod{101} \\
  & 8^8 = 3136 \equiv 5 \pmod{101} \\
  \Rightarrow\quad & c \equiv 8 \cdot 56 \cdot 5 = 2240 \equiv 18 \pmod{101} && (\text{Codewort}) \\
  & a = 18^{77} = 8^{13 \cdot 77} \equiv 8 \pmod{101} && (\text{Klartextwort})
\end{align*}
\end{beispiel}

\subsubsection*{Alice, Bob und Oscar}

Alice verschlüsselt $a \to c = a^e$ und sendet $c$ an Bob. \hfill (Kundin einer Bank)\\
Bob ist im Besitz von $d$ und entschlüsselt $c \to a = c^d$ \hfill (Bank)\\
Oscar versucht den Klartext $a$ zu bekommen.

\begin{beispiel}
Hätte Oscar eine Chance an $a$ zu kommen, wenn er sich $m$ und das verschlüsselte Wort $c$ besorgt?

Ja, ganz einfach. Oscar stellt fest, dass $m = 101$ eine Primzahl ist, dann ist $\varphi(m) = m - 1$. Oscar rechnet selbst $d$ aus
\[
  e\,d \equiv 1 \quad (\bmod\ \varphi(m)) \qquad \Rightarrow \qquad a = c^d
\]
\end{beispiel}

Wie kann das verhindert (oder sehr schwer gemacht) werden?

\begin{definition}[Öffentlicher und privater Schlüssel]
\textbf{Öffentlicher Schlüssel}\\
Alice muss verschlüsseln können
\[
  c = a^e \quad (\bmod\ m)
\]
$\Rightarrow$ \quad öffentlich ist $e$, $m$

\textbf{Privater Schlüssel}\\
Bob muss entschlüsseln können
\[
  a = c^d \quad (\bmod\ m)
\]
Bob hat $d$ aus
\[
  e\,d \equiv 1 \quad (\bmod\ \varphi(m))
\]
berechnet.\\
$\Rightarrow$ \quad privat ist $d$, $\varphi(m)$
\end{definition}

Wie macht es Bob dem Oscar schwer $\varphi(m)$ zu berechnen? Bob setzt
\begin{align*}
  & m = p \cdot q, && p, q \text{ große Primzahlen.} \\
  \Rightarrow\quad & \varphi(m) = (p - 1)(q - 1)
\end{align*}
Aber $p$, $q$ und damit auch $\varphi(m)$ kennt nur Bob.

\begin{beispiel}
Bob wählt $p = 101$, $q = 113$\\
$m = 11413$, \quad $\varphi(m) = 100 \cdot 112 = 11200 = 2^6 \cdot 5^2 \cdot 7$.\\
$e \in \Zm{\varphi(m)}^*$, darf also nicht durch 2, 5 oder 7 teilbar sein. Setzen $e = 3533$.\\
In der Praxis prüft man $e \in \Zm{\varphi(m)}^*$ bei der Berechnung von $d$ mit dem EA.
\begin{align*}
  & d \equiv e^{-1} \pmod{\varphi(m)} \\
  \Rightarrow\quad & d = 6597
\end{align*}
Alice wählt die Zahl $a = 9726 \in \Zm{m}^*$ und berechnet.
\begin{align*}
  c &\equiv a^e = 9726^{3533} \equiv \qquad (\bmod\ m) \\
    &= 5761
\end{align*}
Alice sendet $c$ an Bob und rechnet
\begin{align*}
  a &= 5761^{6597} \qquad (\bmod\ m) \\
    &= 9726
\end{align*}
% UNSICHER: Im Original steht "Alice sendet c an Bob und rechnet"; gemeint ist vermutlich, dass Bob die Entschlüsselung rechnet. Wortlaut unverändert übernommen.
\end{beispiel}

\begin{bemerkung}[Fazit]
Solange es unmöglich ist $m = p \cdot q$ zu faktorisieren, ist RSA ein sicheres Verschlüsselungsverfahren.
\end{bemerkung}
